\documentclass{rvcepaper}
\begin{document}

% ---------- HS351TA Entrepreneurship & IPR (CIE - 1, 09 Oct 2026)
\begin{iempaper}
\paperinfo{dept={INDUSTRIAL ENGINEERING \& MANAGEMENT},date={09\textsuperscript{th} October 2026},exam={CIE -- 1},maxmarks={50},sem={V},prog={UG},duration={90 Min},
 title={ENTREPRENEURSHIP \& INTELLECTUAL PROPERTY RIGHTS},code={HS351TA}}
\iemhead
\begin{cietable}{M}{BT}{CO}
\qpart{Part -- A}
\q{1.}{A software company develops an AI-based tool for legal research and seeks to obtain a patent for the invention in India. Examine the patentability of the AI tool under the Indian Patents Act, 1970, with particular reference to the exclusions relating to computer programmes and software. Discuss the major challenges the company may face in obtaining patent protection.}{10}{4}{5}
\q{2.}{A company has obtained a patent for an innovative electric-vehicle battery management system. A competitor claims that similar technology was disclosed in an earlier patent. At the same time, an investor wants to know the commercial value of the patent. The company is also considering licensing the technology to an automobile manufacturer. Analyse the case and explain how patent searching, valuation, and commercialisation would help the company make an informed decision.}{10}{4}{5}
\q{3.}{A technology company develops an AI-based healthcare application that uses a mathematical algorithm to analyse patient data, presents the results through a graphical interface, and recommends a known medicine for treating a particular disease. The company seeks a patent for the entire system. Analyse the patentability of the invention under the Indian Patents Act, 1970, and identify the relevant categories of non-patentable inventions that may constitute grounds for refusal.}{10}{3}{5}
\q{4.}{``Successful entrepreneurs possess a combination of personal, managerial and behavioural qualities.'' Discuss the important characteristics and traits of successful entrepreneurs, with suitable examples.}{10}{2}{1}
\q{5.}{A technology-based start-up wants to use Artificial Intelligence, digital platforms, social entrepreneurship, sustainable business practices and innovative funding models to establish its business. Explain the emerging trends represented in this case and analyse their impact on entrepreneurship.}{10}{2}{1}
\end{cietable}
\end{iempaper}

% ---------- IM352IA Operations Management (CIE - I A, 09 Oct 2026)
\begin{iempaper}
\paperinfo{dept={INDUSTRIAL ENGINEERING \& MANAGEMENT},date={09\textsuperscript{th} Oct 2026},exam={CIE -- I A},maxmarks={50},sem={V},prog={UG},duration={90 Min},
 title={Operations Management},code={IM352IA}}
\iemhead
\begin{cietable}{M}{BT}{CO}
\q{1.}{With suitable example, Distinguish between Manufacturing and service processes.}{10}{2}{4}
\q{2.}{Identify nine competitive priorities used in operations strategy, and explain how a consistent pattern of decisions can develop organizational capabilities.}{10}{2}{2}
\q{3.}{The owner of a small manufacturing business has patented a new device for washing dishes and cleaning dirty kitchen sinks. Before commercializing the device and adding it to the existing product line, the owner wants to assess its financial feasibility. The estimated variable cost is Rs.7 per unit, while the annual fixed cost is Rs.56,000.
\begin{romi}
\item If the selling price is set at Rs.25 per unit, determine the number of units that must be produced and sold to break even. Use both algebraic and graphic approaches.
\item The forecasted first-year sales are 10,000 units if the selling price is reduced to Rs.15 per unit. Under this pricing strategy, calculate the total contribution to profit during the first year.
\end{romi}}{10}{3}{3}
\q{4.}{A retailer must decide whether to build a small or large facility at a new location. Demand at the location can be either low or high, with probabilities estimated to be 0.4 and 0.6, respectively. If a small facility is built and demand proves to be high, the manager may choose not to expand (payoff = \$223,000) or to expand (payoff = \$270,000). If a small facility is built and demand is low, there is no reason to expand and the payoff is \$200,000. If a large facility is built and demand proves to be low, the choice is to do nothing (\$40,000) or to stimulate demand through local advertising. The response to advertising may be either modest or sizable, with their probabilities estimated to be 0.3 and 0.7, respectively. If it is modest, the payoff is estimated to be only \$20,000; the payoff grows to \$220,000 if the response is sizable. Finally, if a large facility is built and demand turns out to be high, the payoff is \$800,000. Draw a decision tree. Then analyze it to determine the expected payoff for each decision and event node. Which alternative---building a small facility or building a large facility---has the higher expected payoff?}{10}{3}{3}
\q{5.}{A manufacturing company wants to introduce a new product with high variety and low production volume. Explain how it should select an appropriate manufacturing process structure and process strategy.}{10}{3}{4}
\end{cietable}
\end{iempaper}

\end{document}
